\section{Meta reasoning controllers and artifact selection in 2026}
\label{sec:C09:meta2026}
The September 2026 meta-reasoning study separates worker computation from a controller that assesses artifacts, proposes work, evaluates options, and dispatches context from persistent memory. Its analysis distinguishes whether a correct answer appeared, whether the controller recognized it, and whether it selected it. It reports gains alongside overhead that can hurt small-budget settings~\cite{rewrite_metareasoning2026}. The paper's Appendix A explicitly limits component attribution and resource matching. Its shared budget allowance is not proof that every system consumed equal tokens, latency, or dollars. This is a current example of treating inference control as a mechanism rather than equating a larger reasoning transcript with better allocation.

An artifact graph supplies an observable account of the controller's allocation. Nodes represent partial solutions, evidence, checks, and selected outputs; edges record derivation or reuse. Repeated restarts and refinement of verified partial work can have different cost and failure behavior at equal accuracy. Distinguish a copied assertion from independently obtained evidence, and record whether a computation changed the final choice. A verifier invocation that merely confirms the already selected artifact has a different causal role from one that rejects the incumbent and causes a successful replacement. Log both the selection before the check and the selection after it; a final graph without these decision-time snapshots cannot identify that distinction.

An appropriate ablation keeps workers, tools, task instances, and total actual expenditure comparable while varying the control policy. Compare direct control, independent sampling, structured reuse without explicit evaluation, and the full controller. Separately remove persistent artifact retrieval to test whether compact controller context depends on external memory. The resulting experiment can identify which interaction matters; the combined method's headline score cannot assign causal credit to every named stage. The accessible paper specifies the algorithm and prompts; an executable reproduction additionally requires independently established, versioned implementation artifacts.

For stateful deployment, the controller's compact state is itself a memory policy. Its summaries can omit constraints or misstate artifact validity. A worker should receive the evidence needed for its assigned subproblem, and its output should preserve unresolved assumptions. Version changes in workers or tools can invalidate earlier artifacts. The inference controller therefore needs the same compatibility and provenance contracts as other long-running agent components. A successful short benchmark run does not establish indefinite artifact validity or reliable cross-session reuse.



